\chapter{Discussion}

\section{Why Standard Fuzzy TOPSIS Fails}
Classical fuzzy TOPSIS directly aggregates all decision-maker ratings. When coordinated biased decision makers give extreme favorable ratings to a weak target and unfavorable ratings to other alternatives, the aggregated fuzzy matrix becomes distorted. The experiments show that M1 often moves the attacked target to rank 1.

\section{Why Plain Bagging Is Not Enough}
M2 corrects the Random-Forest-style bootstrap logic by sampling decision makers with replacement. However, bagging alone does not identify biased decision makers. If the contaminated group is large enough, many bags still contain biased ratings. Therefore, M2 reduces instability but does not reliably block targeted manipulation.

\section{Why M4 and M6 Are Proposed}
M4 and M6 are both strong because they use reliability. M4 applies reliability to sampling probability, while M6 applies reliability directly to fuzzy aggregation. Their success through 40\% effective contamination shows that reliability-aware methods can protect fuzzy TOPSIS under sub-majority structured attacks.

\section{Why M5 Is Not a Main Proposed Method}
M5 is conceptually attractive because coordinated attackers may form a cluster. However, the experiments show inconsistent behavior. It sometimes helps and sometimes behaves similarly to weaker methods. Therefore, M5 should be discussed as an explored branch or ablation, not as a final proposed method.

\section{Why M7 Is the Final Method}
M7 improves reliability estimation by combining entropy, variance consistency, and clone/agreement-pattern signals. This allows it to detect structured coordinated manipulation more effectively than methods relying only on distance to a median consensus. M7 is the only tested method that preserved the clean target rank across all 18 attack-fraction scenarios.

\section{Threat-Model Boundary}
The method does not detect all possible human bias. Honest human disagreement is normal. Ordinary subjective preference should not automatically be treated as malicious. The limitation arises when malicious decision makers deliberately mimic honest variation. If attackers are statistically indistinguishable from honest decision makers, an unsupervised reliability method cannot guarantee detection without additional information such as historical trust, external ground truth, or supervised labels.

\section{Practical Interpretation}
The proposed framework is best understood as a decision-support safeguard. It does not replace expert judgment. Instead, it alerts the decision process to structurally suspicious rating patterns and reduces their influence in the final ranking.
